% Vietnamese translation for OI Public Library
% Source-File: IOI中国国家候选队论文/2003/许智磊/许智磊.ppt
\documentclass[11pt]{article}
\usepackage{oipl-vn}

\title{Bản trình chiếu: Chuyển sang phần bù trong bài toán thống kê}
\author{Xu Zhilei}
\date{2003}

\begin{document}
\maketitle

\origtitle{Complement transformation in counting problems}
\sourcepath{IOI China candidate team papers / 2003 / Xu Zhilei / Xu Zhilei.ppt}

\section{Tư tưởng của bài trình bày}

Bài trình bày xem bài toán thống kê là bài toán đếm các đối tượng thỏa một
số ràng buộc. Khi đếm trực tiếp quá khó, ta thử đếm toàn bộ \(S\), đếm phần
vi phạm \(T\), rồi lấy đáp án \(R=S-T\). Điểm quan trọng không chỉ là công
thức phần bù, mà là chọn được phía dễ liệt kê hơn.

\section{Ví dụ 1: TRO của POI}

Bài \textit{TRO} cho \(n\) điểm, giữa mọi cặp điểm có một cạnh đỏ hoặc đen,
không có ba điểm thẳng hàng. Dữ liệu cho danh sách cạnh đỏ và yêu cầu đếm số
tam giác đơn sắc. Ràng buộc \(n\le 1000\), \(m\le 250000\) làm cách liệt kê
mọi bộ ba \(O(n^3)\) trở nên quá nặng.

Thử đếm trực tiếp từ một đỉnh thất bại vì hai cạnh cùng màu \(AB,AC\) chưa
đủ quyết định tam giác \(ABC\) đơn sắc; màu của \(BC\) vẫn chưa biết. Phần
bù lại có cấu trúc rõ hơn: một tam giác không đơn sắc luôn có đúng hai cạnh
cùng màu và một cạnh khác màu, nên nó tương ứng với hai cặp cạnh khác màu có
chung đỉnh.

Tổng số tam giác là
\[
  S=\binom n3=\frac{n(n-1)(n-2)}6.
\]
Nếu \(E[i]\) là số cạnh đỏ kề với đỉnh \(i\), thì số cạnh đen kề với \(i\) là
\(n-1-E[i]\). Số cặp cạnh khác màu có chung đỉnh \(i\) là
\[
  E[i](n-1-E[i]).
\]
Do mỗi tam giác không đơn sắc được đếm hai lần theo cách này, đặt
\[
  Q=\sum_i E[i]\,(n-1-E[i])
\]
thì số tam giác đơn sắc là
\[
  R=\frac{n(n-1)(n-2)}6-\frac Q2.
\]
Thuật toán chỉ cần duyệt các cạnh đỏ để tính bậc đỏ của từng đỉnh, rồi cộng
công thức trên, nên đạt \(O(n+m)\) thời gian và \(O(n)\) bộ nhớ.

\section{Ví dụ 2: Sea Battle}

\textit{Sea Battle} cho một bàn \(N\times M\) với \(L\) chiến hạm hình chữ
nhật đã đặt. Cần đếm số cách đặt thêm một hình chữ nhật \(P\times Q\) sao cho
không trùng và không kề tám hướng với các chiến hạm cũ. Dữ liệu có
\(N,M\le 30000\), \(L\le 30\), \(P,Q\le 5\). Liệt kê mọi vị trí trên bàn là
quá lớn, còn rời rạc hóa tuy làm được nhưng phức tạp.

Nếu chưa xét chiến hạm cũ, số vị trí đặt trong một hình chữ nhật
\(X\times Y\) là
\[
  (X-P+1)(Y-Q+1)
\]
khi \(X\ge P\) và \(Y\ge Q\), ngược lại bằng \(0\). Hai hình chữ nhật được
xem là giao hoặc kề nếu các khoảng hàng và khoảng cột của chúng chồng nhau
sau khi mở rộng hình đã đặt thêm một ô về bốn phía:
\[
\begin{aligned}
  &(BX1\le AX2+1)\land(BX2\ge AX1-1)\\
  &\land(BY1\le AY2+1)\land(BY2\ge AY1-1).
\end{aligned}
\]

Chuyển sang phần bù, ta đếm số cách đặt mới giao/kề với ít nhất một chiến
hạm cũ. Với một chiến hạm đã đặt \(A\), mọi vị trí của hình mới \(B\) có thể
giao với \(A\) nằm trong một khung nhỏ quanh \(A\), kích thước không quá
\((2P+X)\times(2Q+Y)\). Vì \(P,Q\) nhỏ và \(L\) nhỏ, lượng liệt kê ở phần bù
rất thấp.

Vấn đề còn lại là đếm lặp: một vị trí của \(B\) có thể giao với nhiều chiến
hạm cũ. Slide xử lý bằng quy tắc thứ tự. Khi xét chiến hạm thứ \(i\), chỉ
cộng một vị trí \(B\) nếu \(B\) giao với chiến hạm \(i\) và không giao với
bất kỳ chiến hạm nào có số thứ tự nhỏ hơn. Như vậy mỗi vị trí vi phạm được
tính đúng một lần, tại chiến hạm đầu tiên mà nó giao.

Độ phức tạp trở thành \(O(P^2Q^2L^2)\): với mỗi chiến hạm, chỉ liệt kê các vị
trí trong khung nhỏ quanh nó, và với mỗi vị trí chỉ kiểm tra giao với tối đa
\(L\) chiến hạm trước đó.

\section{Thông điệp}

Trong \textit{TRO}, chuyển sang phần bù tạo ra một công thức tổ hợp mà phía
thuận không có. Trong \textit{Sea Battle}, phần bù không thay đổi bản chất
đếm, nhưng làm lượng liệt kê giảm từ toàn bàn xuống các vùng rất nhỏ quanh
chiến hạm đã đặt. Hai ví dụ nhấn mạnh cùng một nguyên tắc: nếu đối tượng cần
đếm trực tiếp quá khó, hãy thử tìm phần đối lập có cấu trúc rõ hơn, rồi dùng
thứ tự hoặc công thức tổ hợp để tránh đếm lặp.

\end{document}
